% ==============================================================================
% einstein_approach.tex
% "Sobre la distribución de energía en impactos de nanocúmulos:
%  un análisis en el espíritu de Einstein"
%
% Objetivo pedagógico: mostrar cómo razonar desde primeros principios
% antes de recurrir a la simulación numérica.
% ==============================================================================

\documentclass[11pt, a4paper]{article}

% --- Tipografía y formato ------------------------------------------------------
\usepackage[T1]{fontenc}
\usepackage[utf8]{inputenc}
\usepackage{lmodern}
\usepackage{microtype}
\usepackage[margin=2.8cm, top=3.0cm, bottom=3.2cm]{geometry}
\usepackage{setspace}
\setstretch{1.15}

% --- Matemáticas ---------------------------------------------------------------
\usepackage{amsmath, amssymb, amsthm}
\usepackage{mathtools}

% --- Gráficos y tablas ---------------------------------------------------------
\usepackage{graphicx}
\usepackage[gray]{xcolor}
\usepackage{booktabs}
\usepackage{array}

% --- Encabezado y secciones ---------------------------------------------------
\usepackage{titlesec}
\titleformat{\section}{\large\bfseries}{{\thesection.}}{0.6em}{}
\titleformat{\subsection}{\normalsize\bfseries}{{\thesubsection.}}{0.6em}{}

% --- Cajas para ideas clave ---------------------------------------------------
\usepackage{mdframed}
\mdfdefinestyle{pregunta}{
  linewidth=0.6pt,
  linecolor=black,
  backgroundcolor=white,
  innertopmargin=8pt,
  innerbottommargin=8pt,
  innerleftmargin=12pt,
  innerrightmargin=12pt,
  skipabove=10pt,
  skipbelow=6pt,
}
\mdfdefinestyle{resultado}{
  linewidth=1.2pt,
  linecolor=black,
  backgroundcolor=black!6,
  innertopmargin=8pt,
  innerbottommargin=8pt,
  innerleftmargin=12pt,
  innerrightmargin=12pt,
  skipabove=10pt,
  skipbelow=6pt,
}

% --- Hipervínculos (sin color) ------------------------------------------------
\usepackage[hidelinks]{hyperref}

% --- Datos del documento ------------------------------------------------------
\title{\large\bfseries Sobre la distribución de energía en impactos de nanocúmulos\\[4pt]
       \normalsize\mdseries Un análisis desde primeros principios}
\author{Notas de investigación\thanks{Elaboradas en el contexto del proyecto
        \emph{Distribución de energía en impactos de nanocúmulos LJ},
        inspiradas en el método de razonamiento de A. Einstein.}}
\date{}

% ==============================================================================
\begin{document}
\maketitle
\thispagestyle{empty}

\begin{abstract}
\noindent
Un nanocúmulo de $N$ átomos Lennard-Jones cae sobre una superficie cristalina
bajo gravedad constante. La pregunta experimental es: ¿en qué proporción se
distribuye la energía gravitatoria inicial $E_0 = Mgh_0$ entre el movimiento
del centro de masa, los modos internos del cúmulo y los fonones de la
superficie? Este documento desarrolla el razonamiento analítico completo,
comenzando por el caso elemental $N=1$ y construyendo hacia el caso general.
El objetivo no es reemplazar la simulación sino entenderla.
\end{abstract}

\tableofcontents
\vspace{1em}
\hrule
\vspace{1.5em}

% ==============================================================================
\section{La pregunta en su forma más simple}

Antes de escribir una sola línea de código, Einstein habría formulado el
problema de la manera más desnuda posible. Tomemos esa actitud.

Un objeto cae desde altura $h_0$ y rebota sobre una superficie. Si la
colisión fuera perfectamente elástica, el objeto volvería a $h_0$. No lo hace:
llega sólo a $h_1 < h_0$. La pregunta es \emph{adónde fue la diferencia}
$E_0 - Mgh_1$.

Hay exactamente tres lugares donde puede estar:
\begin{enumerate}
  \item En los \textbf{modos internos del cúmulo} (vibración relativa de los
        átomos de la bola respecto a su centro de masa).
  \item En los \textbf{fonones de la superficie} (oscilaciones de la red
        cristalina excitadas por el impacto).
  \item Ya no está: el cúmulo se \textbf{adhirió} a la superficie
        (sólo si hay atracción bola--superficie, que en este caso no hay).
\end{enumerate}

Definimos las fracciones de energía:
\begin{equation}
  \alpha_\text{CM} = \frac{Mgh_1}{E_0}, \qquad
  \alpha_\text{int} = \frac{\Delta E_\text{bola}}{E_0}, \qquad
  \alpha_\text{surf} = \frac{\Delta E_\text{sup}}{E_0},
\end{equation}
con la restricción de conservación
\begin{equation}
  \alpha_\text{CM} + \alpha_\text{int} + \alpha_\text{surf} = 1.
  \label{eq:conservacion}
\end{equation}

El coeficiente de restitución de alturas $e_1 = \sqrt{h_1/h_0}$ satisface
$\alpha_\text{CM} = e_1^2$, de modo que la energía disipada es $(1-e_1^2)E_0$.

\begin{mdframed}[style=pregunta]
\textbf{La pregunta de Einstein:} ``¿Cuál es el caso más simple que aún
contiene la física esencial?'' Respuesta: $N=1$, un único átomo.
En ese caso $\alpha_\text{int} = 0$ \emph{exactamente} (no hay modos
internos), y toda la disipación debe ir a la superficie.
\end{mdframed}

% ==============================================================================
\section{Experimento mental I: \texorpdfstring{$N=1$}{N=1}}

\subsection{Modelo de colisión elástica con masa efectiva}

Un átomo de masa $m$ con velocidad $v_0 = \sqrt{2gh_0}$ impacta una
superficie. El potencial bola--superficie es \emph{puramente repulsivo}, de
modo que la colisión conserva la energía total. No hay transferencia de
energía química: toda ``pérdida'' de energía cinética del átomo va
\emph{necesariamente} a los fonones de la superficie.

El modelo más simple concebible: la superficie se representa como una
\textbf{masa efectiva} $M^*$ unida a un oscilador (el resto del cristal).
La colisión 1D elástica entre el átomo ($m$, velocidad $v_0$) y la masa
efectiva ($M^*$, en reposo) da:
\begin{align}
  v_\text{átomo}' &= \frac{m - M^*}{m + M^*}\,v_0, \\[4pt]
  V_\text{sup}'   &= \frac{2m}{m + M^*}\,v_0.
\end{align}

La fracción de energía transferida a la superficie es:
\begin{equation}
  \alpha_\text{surf}^{(N=1)} = \frac{M^* (V_\text{sup}')^2}{mv_0^2}
  = \frac{4mM^*}{(m+M^*)^2}.
  \label{eq:alfa_surf_N1}
\end{equation}

El coeficiente de restitución de velocidades es $e_1 = |v_\text{átomo}'|/v_0$,
y puesto que $M^* > m$ (la superficie es más masiva), el átomo rebota y
\begin{equation}
  \boxed{e_1 = \frac{M^* - m}{M^* + m}, \qquad
  \alpha_\text{surf}^{(N=1)} = 1 - e_1^2.}
  \label{eq:e1_N1}
\end{equation}

\noindent Esta es la relación exacta para $N=1$ con potencial puramente
repulsivo. No requiere hipótesis adicionales.

\subsection{La masa efectiva como función de la velocidad}

Despejando $M^*$ de~\eqref{eq:e1_N1}:
\begin{equation}
  M^*(v_0) = m\,\frac{1 + e_1(v_0)}{1 - e_1(v_0)}.
  \label{eq:M_eff}
\end{equation}

\noindent Esta expresión tiene una interpretación directa: $M^*$ es la masa de
la sección de superficie que ``participa'' en la colisión. A baja velocidad,
el tiempo de contacto $\tau_c$ es largo comparado con el período fonónico
$\tau_D = 2\pi/\omega_D$, y muchos átomos de la superficie se mueven
colectivamente~--- $M^*$ es grande. A alta velocidad, $\tau_c$ es corto y
sólo los átomos inmediatamente bajo el impacto responden~--- $M^*$ es pequeño.

De los datos de simulación para $N=1$:
\begin{center}
\begin{tabular}{cccc}
  \toprule
  $v_0$ ($\sigma/\tau$) & $e_1$ & $M^*/m$ & $\alpha_\text{surf}$ \\
  \midrule
  0.74 & 0.979 & $\sim 94$ & 0.041 \\
  0.97 & 0.969 & $\sim 63$ & 0.060 \\
  1.20 & 0.971 & $\sim 66$ & 0.057 \\
  1.46 & 0.946 & $\sim 35$ & 0.105 \\
  1.74 & 0.899 & $\sim 19$ & 0.192 \\
  2.03 & 0.885 & $\sim 16$ & 0.216 \\
  \bottomrule
\end{tabular}
\end{center}

\noindent El resultado central: $M^*(v_0)$ decrece con la velocidad de impacto
porque a mayor $v_0$, el tiempo de contacto $\tau_c \sim \sigma/v_0$ es menor
y menos modos fonónicos de la superficie quedan acoplados al impactador.

\begin{mdframed}[style=resultado]
\textbf{Conclusión para $N=1$:} $\alpha_\text{surf} = 1 - e_1^2$ es exacto.
La dependencia $\alpha_\text{surf}(v_0)$ refleja la \emph{impedancia acústica}
de la superficie: a alta velocidad, la onda de choque excita más fonones de
alta frecuencia y $M^*$ disminuye.
\end{mdframed}

% ==============================================================================
\section{Experimento mental II: \texorpdfstring{$N > 1$}{N mayor que 1} y el conteo de modos}

\subsection{Separación de escalas de tiempo}

Para $N > 1$, la bola tiene modos de oscilación internos con frecuencia
característica $\omega_\text{bola}$. Estimemos esta frecuencia a partir del
potencial LJ.

El mínimo del LJ está en $r_0 = 2^{1/6}\sigma \approx 1.122\,\sigma$.
La constante de fuerza en el mínimo es:
\begin{equation}
  k_\text{LJ} = V''(r_0) = \frac{72\,\varepsilon}{r_0^2}
               \approx \frac{72}{1.260}\,\frac{\varepsilon}{\sigma^2}
               \approx 57.1\,\frac{\varepsilon}{\sigma^2}.
\end{equation}

La frecuencia angular de una vibración entre dos átomos de masa $m=1$:
\begin{equation}
  \omega_\text{bola} = \sqrt{\frac{k_\text{LJ}}{\mu}}
  = \sqrt{\frac{k_\text{LJ}}{m/2}} = \sqrt{2k_\text{LJ}}
  \approx 10.7\,\tau^{-1},
\end{equation}
donde $\mu = m/2$ es la masa reducida. El período de oscilación interno es:
\begin{equation}
  \tau_\text{bola} = \frac{2\pi}{\omega_\text{bola}} \approx 0.59\,\tau.
\end{equation}

El tiempo de contacto con la superficie puede estimarse como:
\begin{equation}
  \tau_c \sim \frac{\sigma}{v_0} \approx \frac{1}{v_0}\,\tau,
\end{equation}
que para $v_0 \in [0.7, 2.0]\,\sigma/\tau$ da $\tau_c \in [0.5, 1.5]\,\tau$.

Comparando las dos escalas:
\begin{equation}
  \frac{\tau_c}{\tau_\text{bola}} \sim \frac{0.5\text{--}1.5}{0.59} \sim 1\text{--}3.
\end{equation}

\begin{mdframed}[style=pregunta]
$\tau_c \gtrsim \tau_\text{bola}$: \textbf{la bola tiene tiempo de redistribuir
energía entre sus modos internos \emph{durante} la colisión}. Esto justifica
tratar la distribución de energía con un argumento de equipartición.
\end{mdframed}

\subsection{Argumento de equipartición de modos}

Durante la colisión, el sistema acoplado \{modos internos de la bola\} $\cup$
\{modos de la superficie en la zona de contacto\} alcanza una distribución de
energía de equilibrio interno. Si cada modo accesible recibe la misma energía
media, la fracción que queda en los modos internos de la bola es simplemente:

\begin{equation}
  \alpha_\text{int} = \frac{f_\text{bola}}{f_\text{bola} + f_\text{sup}},
\end{equation}

donde:
\begin{itemize}
  \item $f_\text{bola} = 3N - 3$ modos internos de la bola
        (descartando las 3 traslaciones del CM),
  \item $f_\text{sup} = 3N_c$ modos de la superficie dentro del área de contacto
        ($N_c$ átomos de superficie efectivos).
\end{itemize}

Sustituyendo:
\begin{equation}
  \boxed{\alpha_\text{int}(N) = \frac{3N - 3}{3N - 3 + 3N_c}
  = \frac{N - 1}{N - 1 + N_c}.}
  \label{eq:gutierrez}
\end{equation}

Esta es la fórmula de Gutiérrez~(2012). Contiene un único parámetro libre,
$N_c$, que cuantifica el acoplamiento bola--superficie.

\subsection{Determinación de \texorpdfstring{$N_c$}{Nc} y verificación}

De la simulación para $N=83$, $\langle\alpha_\text{int}\rangle \approx 0.213$
para $v_0 > 0.9\,\sigma/\tau$. Despejando $N_c$:
\begin{equation}
  N_c = (N-1)\,\frac{1 - \alpha_\text{int}}{\alpha_\text{int}}
      = 82 \times \frac{0.787}{0.213} \approx 303.
\end{equation}

Con $N_c \approx 303$, las predicciones de~\eqref{eq:gutierrez} para
los demás valores de $N$ son:

\begin{center}
\begin{tabular}{ccccc}
  \toprule
  $N$ & Modos $3N-3$ & $\alpha_\text{int}^\text{pred}$
      & $\alpha_\text{int}^\text{sim}$ & discrepancia \\
  \midrule
   1  &    0   &  0.000  &  0.000  & --- \\
   2  &    3   &  0.003  &  0.002  & $-33\%$ \\
   7  &   18   &  0.019  &  0.023  & $+21\%$ \\
  13  &   36   &  0.038  &  0.053  & $+39\%$ \\
  55  &  162   &  0.151  &  0.190  & $+26\%$ \\
  83  &  246   &  0.213  &  0.213  & (ajuste) \\
  \bottomrule
\end{tabular}
\end{center}

La fórmula captura la tendencia: $\alpha_\text{int}$ crece con $N$ porque hay
más modos internos que capturar energía. Las discrepancias del $20$--$40\%$
para $N$ intermedio indican que $N_c$ no es estrictamente constante: el área
de contacto efectiva depende del tamaño del cúmulo.

\subsection{Interpretación geométrica de \texorpdfstring{$N_c$}{Nc}}

¿Qué región de la superficie corresponde a $N_c \approx 303$?

La superficie FCC con parámetro de red $a_0 = 1.55\,\sigma$ tiene densidad
volumétrica $\rho = 4/a_0^3 \approx 1.07\,\text{átomos}/\sigma^3$.
Un hemisferio de radio $R$ bajo el punto de impacto contiene:
\begin{equation}
  N_\text{hem}(R) = \frac{2}{3}\pi R^3\,\rho.
\end{equation}
Igualando $N_\text{hem}(R^*) = 303$ se obtiene $R^* = 5.1\,\sigma$, es decir:
\begin{equation}
  R^* \approx 2\,r_\text{bola} \approx 1.9\,r_\text{bola}.
\end{equation}
Este resultado tiene una lectura física directa: durante el tiempo de
contacto $\tau_c \sim r_\text{bola}/v_0 \approx 1$--$3\,\tau$, la onda
elástica penetra una distancia $d = c_s\,\tau_c$ donde
$c_s \approx 8\,\sigma/\tau$ es la velocidad del sonido LJ. El volumen
excitado es un hemisferio de radio comparable a $d \sim 2r_\text{bola}$,
que contiene precisamente $N_c \approx 300$ átomos de superficie.

\begin{mdframed}[style=resultado]
$N_c \approx 303$ = número de átomos en un hemisferio de radio
$R^* \approx 2\,r_\text{bola} \approx 5\,\sigma$ bajo el punto de impacto.
Esto equivale a la región acústicamente accesible en el tiempo de contacto:
el frente de onda elástica penetra $\sim 3$ distancias interatómicas.
\end{mdframed}

% ==============================================================================
\section{La anomalía: por qué \texorpdfstring{$\Delta V_B > T_B$}{ΔVB > TB}}

\subsection{El resultado inesperado}

La simulación muestra que la energía potencial interna de la bola aumenta
\emph{más} que su energía cinética interna:
\begin{equation}
  \Delta V_B > T_B \quad \text{para } v_0 > 0.9\,\sigma/\tau.
\end{equation}

Para un oscilador \emph{armónico} en equilibrio térmico, el teorema de
equipartición implica $\langle T \rangle = \langle V - V_0 \rangle$. ¿Por qué
la simulación viola esta igualdad?

\subsection{Expansión anarmónica del potencial LJ}

Expandamos el potencial LJ en torno al equilibrio $r_0 = 2^{1/6}\sigma$:
\begin{equation}
  V(r) = V(r_0) + \tfrac{1}{2}k_\text{LJ}(r-r_0)^2
         + \tfrac{1}{6}\gamma(r-r_0)^3 + \cdots
  \label{eq:lj_expansion}
\end{equation}

El coeficiente anarmónico de tercer orden:
\begin{equation}
  \gamma = V'''(r_0) = -\frac{7 \times 72\,\varepsilon}{r_0^3}
         = -\frac{504\,\varepsilon}{r_0^3} < 0.
\end{equation}

El signo negativo de $\gamma$ es fundamental: describe el \textbf{ablandamiento}
del potencial para separaciones $r > r_0$, es decir, los átomos se repelen
menos de lo que un armónico predice cuando se separan.

\subsection{Desplazamiento medio y expansión térmica}

Para un oscilador con perturbación anarmónica, el valor medio de la
elongación puede calcularse por teoría de perturbaciones de primer orden:
\begin{equation}
  \langle r - r_0 \rangle = -\frac{\gamma}{2k_\text{LJ}^2}\,k_B T > 0,
  \label{eq:expansion_termica}
\end{equation}
donde hemos usado $\langle(r-r_0)^2\rangle = k_BT/k_\text{LJ}$
(equipartición del modo armónico) y el teorema de Hellmann-Feynman para
la corrección de primer orden. El resultado positivo confirma la
\textbf{expansión térmica}: los átomos del cúmulo se separan en promedio
cuando vibran.

Sustituyendo los valores numéricos del LJ:
\begin{equation}
  \langle r - r_0 \rangle
  = \frac{504\,\varepsilon/r_0^3}{2\,(57.1\,\varepsilon/\sigma^2)^2}\,k_BT
  = \frac{504}{2 \times 57.1^2}\,\frac{\sigma^2}{r_0^3}\,k_BT
  \approx 0.077\,\frac{k_BT}{\varepsilon}\,\sigma.
\end{equation}

A temperatura interna $k_BT_B \sim \alpha_\text{int} E_0 / (3N-3)$, la
separación media aumenta en $\Delta r \sim 0.077\,k_BT_B/\varepsilon$,
desplazando la PE media hacia valores menos negativos.

\subsection{El exceso de energía potencial}

La energía potencial interna total de la bola es suma sobre todos los pares:
\begin{equation}
  \Delta V_B = \sum_{\langle i,j\rangle}
               \left[\tfrac{1}{2}k_\text{LJ}\langle\delta_{ij}^2\rangle
               + \tfrac{1}{6}\gamma\langle\delta_{ij}^3\rangle\right] - V_0,
\end{equation}
donde $\delta_{ij} = r_{ij} - r_0$. El primer término es $T_B$ (equipartición).
El segundo término es la corrección anarmónica:
\begin{equation}
  \Delta V_B - T_B
  = \frac{\gamma}{6}\sum_{\langle i,j\rangle}\langle\delta_{ij}^3\rangle
  = \frac{\gamma}{6} \cdot N_\text{pares}
    \cdot \left(-\frac{3\gamma}{k_\text{LJ}^3}\right)(k_BT_B)^2.
\end{equation}

Puesto que $\gamma < 0$, este término es positivo: la PE interna supera a
la KE interna en una cantidad proporcional a $T_B^2$. A alta energía de
impacto, el canal potencial $\alpha_{V_B}$ puede dominar sobre $\alpha_{T_B}$.

\begin{mdframed}[style=resultado]
\textbf{Conclusión:} $\Delta V_B > T_B$ no es un error numérico. Es una
consecuencia directa de la anarmónicidad del LJ: el cúmulo, al
calentarse, \emph{se expande} y sus pares quedan en una región donde la
PE media supera a la KE media. Este efecto crece como $T_B^2$ y domina
el canal interno a alta velocidad de impacto.
\end{mdframed}

\subsection{Verificación cuantitativa}

Ajustando $\Delta V_B = a\,T_B + b\,T_B^2$ sobre todos los datos de
$N \in \{2,7,13,55,83\}$ y las seis alturas (filtro $T_B > 0.1\,\varepsilon$):
\begin{equation}
  \Delta V_B = (0.978 \pm 0.153)\,T_B + (0.019 \pm 0.005)\,T_B^2,
  \qquad R^2 = 0.973.
\end{equation}
El coeficiente lineal $a \approx 1$ confirma la equipartición armónica.
El coeficiente cuadrático $b = 0.019\,\varepsilon^{-1}$, positivo con
$3.6\sigma$ de significancia, es la huella de la anarmónicidad LJ.

Una inspección de los coeficientes del Apéndice~A muestra que la
corrección de primer orden en $\gamma$ (término cúbico) y la corrección
de primer orden en $\lambda = V''''(r_0) = 26712\,\varepsilon/r_0^4$
(término cuártico) se \emph{casi} cancelan entre sí:
\begin{equation}
  c_\text{cúbico} = -\frac{5\gamma^2}{12k^3} \approx -\frac{0.638}{\varepsilon},
  \qquad
  c_\text{cuártico} = \frac{\lambda}{8k^2} \approx +\frac{0.644}{\varepsilon}.
\end{equation}
El cuártico domina por apenas el $1\%$, de modo que el neto
$c_\text{neto} \approx +0.006\,\varepsilon^{-1}$ tiene el signo correcto
(positivo, $\Delta V_B > T_B$) pero su magnitud requiere términos de orden
superior: en nuestro régimen $k_BT_\text{ball}/\varepsilon \sim 0.2$,
la perturbación de primer orden subestima $b$ por un factor $\sim 3$.
La concordancia de signo es suficiente para la interpretación física.

% ==============================================================================
\section{Los tiempos de relajación}

Einstein no se habría contentado con conocer las fracciones de energía al
primer rebote. Habría preguntado: \textbf{¿en cuánto tiempo se equilibra el
sistema completo?}

\subsection{Tres escalas de tiempo}

Identificamos tres tiempos característicos:

\paragraph{$\tau_c$: tiempo de contacto.}
Escala de la interacción bola--superficie:
\begin{equation}
  \tau_c \sim \frac{\sigma}{v_0} \approx 0.5\text{--}1.5\,\tau.
\end{equation}
Es el tiempo durante el cual ocurre la transferencia de energía.

\paragraph{$\tau_\text{int}$: tiempo de equilibración interna de la bola.}
Tiempo para que los modos internos se termalicen entre sí. En un cristal
finito de $N$ átomos y parámetro de red $a_0$, con velocidad del sonido
$c_s \sim a_0\,\omega_\text{bola} \approx 1.5\,\sigma/\tau$:
\begin{equation}
  \tau_\text{int} \sim \frac{N^{1/3} a_0}{c_s} \approx N^{1/3} \times 1\,\tau.
\end{equation}
Para $N=13$: $\tau_\text{int} \approx 2.3\,\tau$. Para $N=83$:
$\tau_\text{int} \approx 4.4\,\tau$. Ambos son comparables a $\tau_c$, lo
que justifica el argumento de equipartición de modos.

\paragraph{$\tau_\text{eq}$: tiempo de equilibración bola--superficie.}
Este es el tiempo para que la bola transfiera toda su energía interna a la
superficie a través de rebotes sucesivos. Cada rebote transfiere una fracción
$\alpha_\text{surf} \approx 0.10$ de la energía interna. El número de rebotes
necesarios es $n^* \sim 1/\alpha_\text{surf} \approx 10$, y el período de
cada rebote es:
\begin{equation}
  \tau_\text{rebote} = 2\sqrt{\frac{2h^*}{g}} \approx 60\,\tau,
\end{equation}
donde $h^* \sim T_B/(Mg)$ es la altura de saturación predicha por
Gutiérrez~(2012). En total:
\begin{equation}
  \tau_\text{eq} \sim \frac{\tau_\text{rebote}}{\alpha_\text{surf}}
  \approx \frac{60}{0.10}\,\tau = 600\,\tau.
\end{equation}

Nuestra simulación corre $400\,\tau$: \textbf{el sistema no ha alcanzado el
equilibrio}. Para observar la equipartición completa bola--superficie
necesitamos correr al menos $5\,\tau_\text{eq} \approx 3000\,\tau$.

\subsection{Separación de escalas}

El cuadro completo de escalas de tiempo es:
\begin{equation}
  \underbrace{\tau_c \sim 1\,\tau}_{\text{contacto}}
  \;\ll\;
  \underbrace{\tau_\text{rebote} \sim 60\,\tau}_{\text{periodo oscilación}}
  \;\ll\;
  \underbrace{\tau_\text{eq} \sim 600\,\tau}_{\text{equilibrio térmico}}.
  \label{eq:escalas}
\end{equation}

Esta separación de escalas \emph{explica} por qué podemos medir una
partición de energía bien definida al primer apogeo: la bola se ha
``termalizado'' internamente ($\tau_c > \tau_\text{int}$) pero no con la
superficie ($\tau_\text{eq} \gg \tau_\text{rebote}$).

\begin{mdframed}[style=resultado]
\textbf{El cuadro físico completo:} (i) Durante el contacto ($\tau_c$),
la bola distribuye energía entre sus modos internos y los modos de
superficie accesibles. (ii) Entre rebotes, la bola es un reservorio de
energía interna casi aislado. (iii) Rebote a rebote, la bola pierde
energía a la superficie lentamente ($\alpha_\text{surf}$ por rebote).
(iv) En tiempo $\tau_\text{eq}$, el sistema completo se equilibra.
\end{mdframed}

% ==============================================================================
\section{Conducta límite de la energía interna}

\subsection{El estado estacionario predicho}

Gutiérrez~(2012) predice que, después de muchos rebotes, la energía interna
de la bola converge a un valor límite $T_B^*$ tal que la bola ya no puede
llegar más alto que:
\begin{equation}
  h^* = \frac{T_B^*}{Mg}.
  \label{eq:h_star}
\end{equation}

Podemos estimar $T_B^*$ de nuestro argumento: en cada rebote, la bola
absorbe $\alpha_\text{int}(1-\alpha_\text{CM})E_n$ de energía interna y
pierde $\alpha_\text{surf}T_B^{(n)}$ a la superficie. En estado
estacionario estas dos cantidades se igualan:
\begin{equation}
  \alpha_\text{int}(1-\alpha_\text{CM}^*)E^* = \alpha_\text{surf}\,T_B^*,
\end{equation}
donde $E^* = Mgh^*$. Resolviendo:
\begin{equation}
  T_B^* = \frac{\alpha_\text{int}(1-\alpha_\text{CM}^*)}
               {\alpha_\text{surf}}\,Mgh^*.
\end{equation}

Dado que $T_B^* = Mgh^*$ por~\eqref{eq:h_star}, se obtiene la
consistencia:
\begin{equation}
  1 = \frac{\alpha_\text{int}(1-\alpha_\text{CM}^*)}{\alpha_\text{surf}},
\end{equation}
que fija la relación entre los tres canales en el estado estacionario.
Con los valores medidos ($\alpha_\text{int} \approx 0.21$,
$\alpha_\text{surf} \approx 0.10$), esto predice
$\alpha_\text{CM}^* \approx 0.52$, es decir, la bola en estado
estacionario rebota a aproximadamente la mitad de su altura.

\subsection{Estimación de \texorpdfstring{$h^*$}{h*}}

Usando la conservación de energía en el estado estacionario:
\begin{align}
  Mgh_0 &= E_\text{total inicial}, \\
  T_B^* + Mgh^* &= \text{energía disponible tras equilibración parcial}.
\end{align}

Un balance más cuidadoso (que requiere la simulación larga del Paso~4)
predice $h^* \approx \alpha_\text{int} h_0 / (1 - \alpha_\text{CM}^*)$.
Para $h_0 = 28\,\sigma$: $h^* \approx 12$--$14\,\sigma$.

% ==============================================================================
\section{Predicciones verificables sin simulación adicional}

Reunimos las predicciones cuantitativas del análisis, ordenadas por
jerarquía de certeza:

\begin{enumerate}
  \item \textbf{Exacto (sin parámetros libres):}
        Para $N=1$, $\alpha_\text{surf} = 1 - e_1^2$ en todo rango de $v_0$.

  \item \textbf{Un parámetro ($N_c$):}
        $\alpha_\text{int}(N) = (N-1)/(N-1+N_c)$, con $N_c \approx 303$.
        La predicción puede fallar si $N_c$ depende de $N$ (contacto
        no escala).

  \item \textbf{Cualitativo:}
        $\Delta V_B > T_B$ para $v_0 > 0.9\,\sigma/\tau$, y el exceso
        crece como $T_B^2$ (verificable ajustando $\Delta V_B$ vs $T_B$).

  \item \textbf{Predicción de $\tau_\text{eq}$:}
        El sistema no equilibra bola--superficie antes de $\sim 600\,\tau$.
        Para verificarlo, correr $5000\,\tau$ (Paso~4 del plan).

  \item \textbf{Adhesión ($c_{ps}=1$):}
        Cuando la atracción bola--superficie es no nula, existe un
        umbral $h_0^*(N)$ por debajo del cual la bola no rebota.
        La escala de energía relevante es la energía de adhesión
        $U_\text{ad} \approx N_c \varepsilon$, que predice
        $h_0^* \sim N_c\varepsilon / (Mg)$.
\end{enumerate}

% ==============================================================================
\section{La lección metodológica}

Einstein habría llegado a estos resultados \emph{antes} de encender
la computadora. El método tiene una estructura invariante:

\begin{enumerate}
  \item \textbf{Caso más simple primero.} $N=1$ es analíticamente exacto.
        El resto del problema se construye sobre ese caso límite.

  \item \textbf{Separación de escalas.} Antes de cualquier cálculo,
        comparar $\tau_c$, $\tau_\text{int}$, $\tau_\text{eq}$.
        Las aproximaciones válidas dependen de esas comparaciones.

  \item \textbf{Argumentos de conteo.} La fórmula~\eqref{eq:gutierrez}
        no requiere resolver la dinámica: sólo contar modos y aplicar
        equipartición. Esto es lo que hace Einstein siempre: busca el
        argumento más general posible.

  \item \textbf{Identificar el parámetro libre.} $N_c$ no se puede
        predecir sin información de la geometría del contacto. La
        simulación sirve para medir ese único parámetro.

  \item \textbf{La anomalía como punto de partida.} $\Delta V_B > T_B$
        no es un resultado a ignorar: es la señal de que la física es
        más rica que el modelo armónico. Einstein siempre persiguió las
        discrepancias, no los acuerdos.
\end{enumerate}

\begin{mdframed}[style=resultado]
\textbf{El rol de la simulación en este esquema:} no es explorar, sino
\emph{medir parámetros de un modelo ya comprendido} ($N_c$, $\tau_\text{eq}$,
el prefactor anarmónico) y \emph{verificar predicciones analíticas} en los
casos donde el análisis es exacto ($N=1$). La simulación que no tiene
un modelo analítico al lado que interpretar produce números, no comprensión.
\end{mdframed}

% ==============================================================================
\appendix
\section{Derivación del coeficiente anarmónico \texorpdfstring{$\gamma$}{γ}}

El potencial LJ es $V(r) = 4\varepsilon[(\sigma/r)^{12} - (\sigma/r)^6]$.
Calculamos las derivadas en el mínimo $r_0 = 2^{1/6}\sigma$:
\begin{align}
  V'(r)   &= 4\varepsilon\left[-12\,\frac{\sigma^{12}}{r^{13}}
              + 6\,\frac{\sigma^6}{r^7}\right], \\
  V''(r)  &= 4\varepsilon\left[156\,\frac{\sigma^{12}}{r^{14}}
              - 42\,\frac{\sigma^6}{r^8}\right], \\
  V'''(r) &= 4\varepsilon\left[-2184\,\frac{\sigma^{12}}{r^{15}}
              + 336\,\frac{\sigma^6}{r^9}\right].
\end{align}

Evaluando en $r_0 = 2^{1/6}\sigma$ (donde $\sigma^6/r_0^6 = 1/2$):
\begin{align}
  k_\text{LJ} = V''(r_0) &= 4\varepsilon\left[\frac{156}{r_0^2}\cdot\frac{1}{4}
    - \frac{42}{r_0^2}\cdot\frac{1}{2}\right]
    = 4\varepsilon\,\frac{156/4 - 42/2}{r_0^2}
    = 4\varepsilon\,\frac{39 - 21}{r_0^2}
    = \frac{72\varepsilon}{r_0^2}, \\[6pt]
  \gamma = V'''(r_0) &= 4\varepsilon\left[-\frac{2184}{r_0^3}\cdot\frac{1}{8}
    + \frac{336}{r_0^3}\cdot\frac{1}{4}\right]
    = 4\varepsilon\,\frac{-273 + 84}{r_0^3}
    = -\frac{756\varepsilon}{r_0^3}.
\end{align}

El cociente de anarmónicidad:
\begin{equation}
  \frac{\gamma^2}{k_\text{LJ}^3}
  = \frac{(756)^2\varepsilon^2/r_0^6}{(72)^3\varepsilon^3/r_0^6}
  = \frac{571536}{373248\,\varepsilon}
  \approx \frac{1.53}{\varepsilon}.
\end{equation}

La corrección anarmónica a la energía potencial media por par (a temperatura
$T$ del modo) es de orden $(k_BT)^2/\varepsilon$, confirmando que el exceso
$\Delta V_B - T_B$ es no despreciable cuando $k_BT_B \sim 0.1\text{--}0.3\,\varepsilon$.

\end{document}
